\section{Conclusion}

This work investigated machine-learning-based classification of eight power quality disturbances (PQDs) from the normal class, using eight features extracted from the S transformed signal. The feature-selection analysis showed that the relevance of a feature depends on the disturbance class. The recommended subsets were generally consistent with the expected physical relationships and with the graphical and statistical analyses. In particular, harmonic-content features, voltage-amplitude statistics, and high-frequency features were repeatedly identified as informative. The results also showed that some classes can be separated using only a small number of features, whereas more subtle disturbance classes, such as Flicker, benefit from combinations of more features and nonlinear classifier architectures.

For the multiclass problem, all eight features were retained because every feature was needed to distinguish at least one class. Decision tree, random forest, and multilayer perceptron models were optimized using stratified five-fold cross-validation. On the held-out test set, the random forest was the strongest individual model, achieving a macro F1-score of 0.9822 and a normal-class precision of 0.9281. Combining the three classifiers with soft voting improved both measures: the bagged model achieved a macro F1-score of 0.9844 and a normal-class precision of 0.9412. This improvement indicates that the individual models make partly complementary errors and that their combination is beneficial for robust multiclass classification.

The confusion matrices show that the main remaining ambiguity is between the Normal, Transient, and Flicker classes. This limitation is consistent with the feature distributions, which overlap substantially for these classes. Furthermore, the evaluation is based on the available labelled dataset and a single held-out test split. Consequently, the reported performance should not be interpreted as sufficient evidence for direct deployment in a power-system protection or monitoring application. Additional evaluation on measurements from different operating conditions and networks would be required to assess generalization and reliability.

Future work should therefore focus on collecting and validating more diverse real-world data, improving the separation of Normal, Transient, and Flicker signals, and evaluating confidence-aware or cost-sensitive decision rules that explicitly penalize undetected disturbances. Further improvements could also be obtained by testing additional signal features, calibration methods, and ensemble strategies. Despite these limitations, the results demonstrate that the selected features and the ensemble approach provide a strong foundation for automated PQD classification.
